\section{Price Dynamics and Preference Updating}
\label{sec:dynamics}

\subsection{Price and Shadow Price Relationship}

Following \citet{lange1936economic}, shadow prices are proportional to
market prices.
This relationship also follows directly from
equations~\eqref{eq:demand} and~\eqref{eq:shadow_prices}.

At the CRRA optimum, the first-order condition
$\alpha^{agg}_{i,t}\,C_i^{-\sigma_i} = \mu_t P_{i,t}$
(see equation~\eqref{eq:crra_demand}) implies that the shadow price of
good $i$ satisfies $\pi_i = \mu_t P_{i,t}$, so that:
   
\begin{equation}\label{eq:price_proportional}
    \vec{P}_t = \frac{1}{\mu_t}\,\vec{\pi}_t.
\end{equation}
 
The marginal utility of income $\mu_t$ thus acts as the common
proportionality factor linking shadow prices to market prices. Its
value is pinned down by the budget constraint
$\vec{P}_t^\top \vec{C}_t = Y_t$: substituting
equation~\eqref{eq:demand} gives

\begin{multline}\label{eq:mu_determination}
    Y_t = \sum_{i=1}^n P_{i,t}
          \left(\frac{\alpha^{agg}_{i,t}}{\mu_t P_{i,t}}\right)^{1/\sigma_i}
        = \\ \sum_{i=1}^n \frac{(\alpha^{agg}_{i,t})^{1/\sigma_i}}{\mu_t^{1/\sigma_{i}}P_{i,t}^{1/\sigma_i - 1}} \\ =
       \frac{1}{\mu_t} \sum_{i=1}^n \frac{(\alpha^{agg}_{i,t})^{1/\sigma_i}}{\pi_{i,t}^{1/\sigma_i - 1}}
\end{multline}

which implicitly defines $\mu_t$ as a function of prices, income, and
preference weights. Similarly the Gross domestic product or the total factor income is given by the following:

\begin{multline}
    \text{GDP}_{t} = \frac{1}{\mu_{t}} \vec{\pi}_{t}^{\top} \vec{F}_{t} \\ =
    \frac{1}{\mu_{t}}\left(\lambda_{L,t}\vec{l} + B^{\top} \vec{\lambda}_{K,t} \right)^{\top} \vec{X}_{t} = \vec{v}_{t}^{\top} \vec{X}_{t}
\end{multline}

Assuming the price level stays constant throughout the simulation we define the layspers price index.

\begin{multline}
    \frac{1}{\mu_{t}} \vec{\pi}_{t} \cdot \vec{C}_{t-1} =  \frac{1}{\mu_{t-1}} \vec{\pi}_{t-1} \cdot \vec{C}_{t-1} \\ \implies
    \frac{\mu_{t}}{\mu_{t-1}} = \frac{\vec{\pi}_{t} \cdot \vec{C}_{t-1}}{\vec{\pi}_{t-1} \cdot \vec{C}_{t-1}}
\end{multline}

As the initial income numeraie is known the initial marginal utility of income, the income may be calculated into the future, with the use of the marginal utility of income, and similarly the nominal prices may be calculated by the chained layspers index.

\subsection{Price Adjustment} \label{stability}
In each sub-period $\tau$ within the quarterly loop:
\begin{equation}\label{eq:tatonnement}
    P_{i,t,\tau}^{(k+1)} = P^{(k)}_{i,t}
    \left(1 - \frac{1}{\epsilon_i}
              \cdot \frac{Z^{(k)}_{i,t,\tau}}{C_{i,t,\tau}}
    \right),
\end{equation}
where $Z_{i,t,\tau} = C_{i,t,\tau}(\vec{P}_{t}) -
C^{\mathrm{m}}_{i,t,\tau}$ is excess demand and $\epsilon_i = -1/\sigma_i < 0$
is the price elasticity implied by the CRRA demand
system~\eqref{eq:demand}. The price adjustment rule applies a single
Newton step on the inverse demand function, using good-specific
curvature parameters $\sigma_i$ to calibrate the step size against
observed sectoral excess demand.
